\section{Day 41: Fundamental Theorem of Galois Theory (Mar.\ 11, 2026)}
Recall the fundamental theorem of Galois theory,
\begin{theorem}
    Given $E/F$ a splitting extension of a field $F$ with characteristic zero, there is a bijection
    \[ \{K : E/K/F\} \leftrightarrow \{H < G \coloneq \Gal(E/F)\} \]
    where $K \mapsto \Gal(E/K)$ and $H \mapsto E_H$.
\end{theorem}
In addition, \begin{parlist} \item the bijection reverses inclusions, i.e., if $K_1 \subset K_2$, then $\Gal(E/K_2) < \Gal(E/K_1)$, and if $H_1 < H_2$, we have $E_{H_2} \subset E_{H_1}$. \item We also have $\coor{E:K} = \abs{\Gal(E/K)}$, and $\coor{K:F} = \coor{G:\Gal(E/K)}$, and that \item $\Gal(E/K) < \Gal(E/F)$ being normal (denoted respectively as $H < G$) is equivalent to $E_H$ being a splitting field of $F$, i.e., $\Gal(E_H/F) = G/H$. \end{parlist} As an example,
\[ E = \QQ(x^3 - 2) = \QQ(2^{1/3}, \rho \coloneq e^{2\pi i /3}) = \QQ(2^{1/3}, 2^{1/3} \rho, 2^{1/3} \rho^2), \]
which we once again denote as $\QQ(u_1, u_2, u_3)$. We have the corresponding diagrams
\[ \begin{tikzcd}
    & E = \QQ(u_1, u_2, u_3) & \\
    \QQ(x^2 + x + 1) = \QQ(\rho) \arrow[ur, "3"] & & \QQ(u_3) \arrow[ul, "2"] \\ & \QQ \arrow[ul, "2"] \arrow[ur, "3"]
\end{tikzcd} \]
with the opposite direction for groups, i.e.,
\[ \begin{tikzcd}
    & \{e\} \arrow[dl, "3"] \arrow[dr, "2"] & \\
    \left<(1 \, 2 \, 3)\right> \arrow[dr, "2"] & & \left<(1 \, 2)\right> \arrow[dl, "3"] \\
    & \Gal(E/\QQ) \cong S_3.
\end{tikzcd} \]
Note that $(1 \, 2)$ corresponds to the automorphism swapping $u_1$ and $u_2$ as before, with the same for $(2 \, 3)$ and $(1 \, 3)$, which have corresponding fields $\QQ(u_1)$, $\QQ(u_2)$. In particular, observe that
\[ \rho = \frac{u_{i+1}}{u_i} \]
for all $i$, so $(1 \, 2 \, 3)$ and $(1 \, 3 \, 2)$ do not change $\rho$. Here, everything matches up as in the theorem. As an example, consider $E = \QQ(3x^5 - 15x + 5)$, for which the polynomial on the inside (which we will call $P$) is irreducible by Eisenstein's criterion. Let $G \coloneq \Gal(E/\RR)$.
\begin{lemma}
    In a characteristic zero field, if $P$ is irreducible and $\deg P = n$, then, in splitting fields of $P$, $P$ has exactly $n$ roots.
\end{lemma}
\begin{proof}
    We start by showing that $\gcd(P, P') = 1$. If $\deg P' < \deg P$ with $P' \neq 0$, since we are in a characteristic zero field and the only thing dividing $P$ with degree less than $\deg P$ is $1$, hence $P$ has $5$ roots. If $u$ is a root of $P$, then $[\QQ(u) : \QQ] = \deg P = 5$. Thus, $5 \mid [E : \QQ]$, implying $5 \mid \abs{G}$ by the fundamental theorem of Galois theory. As before, $G < S_5$ as its elements must permute the roots of $P$, otherwise it fixes everything. Thus, $G$ must contain a $5$ cycle, although we don't know which one. We may guess from elementary analysis that
    \[ P' = 15x^4 - 15 = 15(x^4 - 1), \]
    so $P' = 0$ at $x = \pm 1$; we have $P(1) = -7$, $P(-1) = 17$, $P(0) = 5$, so there are exactly $3$ real roots and $2$ complex roots. Complex conjugation must interchange the complex roots, so there is a $2$-cycle in $G$. There are two possibilities: either $G$ is generated by $(1 \, 2 \, 3 \, 4 \, 5)$ and $(1 \, 2)$, or it is generated by $(1 \, 2 \, 3 \, 4 \, 1)$ and $(1 \, 3)$; in either case, it is isomorphic to $S_5$.
\end{proof}
To prove the theorem, we will first show $K = E_{\Gal(E/K)}$.
\begin{proof}
    We already have that $K$ is included in $E_{\Gal(E/K)}$ by definition; the other direction is harder. Suppose $v \in E \setminus K$; we want to show that there exists $\phi \in \Gal(E/K)$ such that $\phi(v) \neq v$. Let $P$ be the minimum polynomial of $v$ over $K$, and $\deg P > 1$ (as $v \notin K$). $P$ has no multiple roots, so $P$ has a root $u \neq v$, and $u \in E$ as $E$ is splitting in the theorem. Thus, $\phi_0$ exists as a minimum polynomial, so is irreducible, with $\phi_0(v) = u$ and $\restr{\phi_0}{K} = \id$. $\phi$ is an extension of $\phi_0$ in $\Aut(E)$ with $\phi(v) = u \neq v$ and $\restr{\phi}{K} = \id$. Thus, we have found the desired map.
\end{proof}
